\chapter{Peta Perjalanan}
\label{ch:peta}

Sebelum masuk ke materi, ada baiknya melihat peta besarnya dulu: mata kuliah apa saja yang
dilalui, dalam urutan apa, dan bagaimana semuanya tersusun menjadi buku ini.

\section{Empat semester}

\Cref{tab:perjalanan} menyusun mata kuliah menurut semester ketika saya mengikutinya. Urutan
bab di buku tidak persis mengikuti urutan waktu. Bab disusun menurut alur gagasan: fondasi
lebih dulu, lalu deep learning, lalu simulasi. Tabel ini berguna bagi pembaca yang ingin
membaca buku ini sambil mengikuti kuliah serupa.

\begin{table}[ht]
\centering\small
\caption{Mata kuliah menurut semester dan bab yang membahasnya.}
\label{tab:perjalanan}
\begin{tabularx}{\linewidth}{@{}p{3.6cm}Xp{1.4cm}@{}}
\toprule
Semester & Mata kuliah & Bab \\
\midrule
Musim dingin 2024/25
 & Deep Learning and Neural Nets I & \ref{ch:dlbasic} \\
 & LSTM and Recurrent Neural Nets & \ref{ch:lstm} \\
 & Computer Vision & \ref{ch:cv} \\
 & Machine Learning: Supervised Techniques & \ref{ch:data} \\
 & Control Systems & \ref{ch:kendali} \\
 & Production Automation Systems & \ref{ch:kendali} \\
 & Optimum and Adaptive Signal Processing Systems & \ref{ch:sinyal} \\
 & AI in Society, AI and Law I & \ref{ch:manusia} \\
 & Patent and Literature Searching & \ref{lamp:riset} \\
\midrule
Musim panas 2025
 & Deep Learning and Neural Nets II & \ref{ch:arsitektur}, \ref{ch:praktik} \\
 & Deep Reinforcement Learning & \ref{ch:rl} \\
 & Machine Learning: Unsupervised Techniques & \ref{ch:unsup} \\
 & Machine Learning and Pattern Classification & \ref{ch:unsup} \\
 & Statistical Learning Theory & \ref{ch:teori} \\
 & Introduction to Autonomous Vehicles & \ref{ch:robot} \\
 & Introduction to Robotic Systems & \ref{ch:robot} \\
 & Learning from User-generated Data & \ref{ch:rekomendasi} \\
 & Introduction to Genetics; Genome Analysis and Transcriptomics & \ref{ch:genomik} \\
 & Structural Bioinformatics & \ref{ch:protein} \\
 & AI in Life Sciences & \ref{ch:hayati} \\
 & Deep Learning for Medical Imaging & \ref{ch:medis} \\
 & Robopsychology; Technology and Society; AI and Law II & \ref{ch:manusia} \\
 & Quantum Information & \ref{ch:kuantum} \\
\midrule
Musim dingin 2025/26
 & Machine Learning: Advanced Techniques & \ref{ch:mlat} \\
 & Probabilistic Models & \ref{ch:prob} \\
 & AI and Visualization & \ref{ch:xai} \\
 & Machine Learning: Basic Techniques & \ref{ch:data}, \ref{ch:stat} \\
 & Reinforcement Learning & \ref{ch:rl} \\
 & Numerical Methods in Fluid Mechanics & \ref{ch:fluida}, \ref{ch:numerik} \\
 & Computational Physics I & \ref{ch:fiskom} \\
 & Optimum and Adaptive Signal Processing Systems (latihan) & \ref{ch:sinyal} \\
 & Radar System Engineering (latihan) & \ref{ch:sinyal} \\
 & Pervasive Computing: Systems and Environments; Design and Development & \ref{ch:pervasif} \\
 & KS Operations Research & \ref{ch:klasik} \\
 & Natural Language Processing & \ref{ch:nlp} \\
 & Communicating AI & \ref{ch:manusia} \\
\midrule
Musim panas 2026
 & Introduction to AI & \ref{ch:klasik} \\
 & Deep Learning: Geometric Techniques & \ref{ch:gdl} \\
 & Numerical Optimization for AI & \ref{ch:optimasi} \\
 & Tugas akhir: recurrence CFD & \ref{ch:rcfd} \\
\bottomrule
\end{tabularx}
\end{table}

Bab-bab tentang \emph{scientific machine learning} (\cref{ch:pinn,ch:operator,ch:penemuan,ch:hibrida})
tidak berasal dari satu mata kuliah. Isinya dirangkai dari beberapa kuliah, terutama machine
learning lanjut, metode numerik fluida, dan fisika komputasi, ditambah bacaan selama
mengerjakan tugas akhir.

Program master ini mensyaratkan 120 ECTS untuk lulus, dan tabel di atas memuat lebih banyak dari itu. Mata kuliah yang melebihi
syarat kelulusan tetap masuk ke buku, sebab merekalah yang membuat ceritanya utuh: robot yang bergerak, molekul yang dirancang dengan
bantuan mesin, sinyal radar yang harus dipisahkan dari derau, rekomendasi yang kita terima setiap hari, dan pertanyaan hukum serta
psikologis yang muncul ketika semua itu bertemu manusia.

Buku ini terbagi menjadi enam bagian. Bagian I berisi fondasi: apa artinya belajar dari data, bagaimana statistik dan Bayes membantu
memilih model, apa yang bisa dijamin teori belajar, bagaimana mencari parameter terbaik, bagaimana AI klasik mencari dan merencanakan,
dan bagaimana menemukan struktur tanpa label. Bagian II membahas deep learning dan persepsi: teknik dasar dan praktiknya, ingatan dari
RNN sampai xLSTM, bahasa, penglihatan, model generatif, simetri, penalaran probabilistik, model difusi, reinforcement learning, dan
explainable AI. Bagian III turun ke dunia sinyal dan sistem fisik: estimasi dan filter, radar, kendali, komputasi pervasif, dan robot.
Bagian IV masuk ke ilmu hayati: genom, molekul dan obat, protein, dan citra medis. Bagian V membahas AI untuk simulasi fisika, bidang
tugas akhir saya. Bagian VI menutup dengan sistem rekomendasi, hukum dan manusia, keselamatan AI, informasi kuantum, dan perkembangan AI
terbaru. Tiga lampiran melengkapinya: glosarium, bekal matematika, dan panduan singkat untuk meneliti, menerbitkan, dan mematenkan.

Hampir setiap bab punya tiga lapis. Lapis pertama adalah isi kuliah, ditulis ulang dengan penurunan dan contoh hitungan. Lapis kedua
adalah bagian ``di luar kelas'' yang menghubungkan materi dengan bidang lain di buku ini. Lapis ketiga adalah bagian \emph{Perkembangan
riset}, yang merangkum makalah-makalah beberapa tahun terakhir. Semua makalah di lapis ketiga diperiksa terhadap catatan resminya, dan
isinya ditulis berdasarkan abstrak yang benar-benar dibaca, bukan dari ingatan.

\section{Cara membaca buku ini}

Buku ini bisa dibaca dengan dua cara, dan keduanya sama sahnya.

Pembaca yang ingin mengikuti setiap penurunan akan terbantu oleh tiga bekal: kalkulus, terutama turunan dan aturan rantai; aljabar
linear, terutama perkalian matriks, nilai eigen, dan dekomposisi nilai singular; serta peluang dan statistik dasar, termasuk
ekspektasi, varians, distribusi Gauss, dan teorema Bayes. Semuanya dirangkum di \cref{lamp:matematika}, yang bisa dibaca lebih dulu
atau dibuka setiap kali sebuah rumus terasa asing. Rumus di buku ini jarang muncul tanpa diturunkan, dan contoh hitungannya sengaja
dibuat kecil supaya bisa diperiksa dengan pensil.

Pembaca yang mencari wawasan tidak perlu berhenti di setiap rumus. Paragraf di sekitar persamaan selalu menjelaskan maknanya dengan
kata-kata, sehingga persamaan bisa dilewati tanpa kehilangan alur. Bab-bab tentang AI dan manusia, keselamatan AI, dan perkembangan
terbaru (\cref{ch:manusia,ch:keselamatan,ch:terbaru}) hampir tidak membutuhkan matematika sama sekali, dan bisa langsung dibaca kapan
saja. Begitu juga bab-bab penerapan, seperti ilmu hayati atau robot, yang bisa dibaca sebagai cerita tentang apa yang dilakukan AI di
bidang-bidang itu.

Bagi yang ingin membaca sebagian saja, \cref{fig:jalur} menyarankan lima jalur. Setiap jalur bisa diikuti sendiri, dan bab-bab yang
dilewati selalu bisa dikunjungi belakangan lewat rujukan silang di dalam teks dan indeks di akhir buku.

\begin{figure}[ht]
\centering
\begin{tikzpicture}[x=1.05cm,y=0.95cm,
  bab/.style={draw=abu,rounded corners=2pt,minimum width=0.85cm,minimum height=0.55cm,inner sep=1pt,font=\footnotesize},
  jalur/.style={font=\footnotesize\itshape,anchor=east}]
  \newcommand{\baris}[3]{%
    \node[jalur] at (-0.6,#1) {#2};
    \foreach \lab [count=\i from 0] in {#3} {
      \node[bab,fill=biru!12] (b#1-\i) at (\i*1.05,#1) {\ref{\lab}};
      \ifnum\i>0 \pgfmathtruncatemacro{\j}{\i-1}\draw[->,abu] (b#1-\j) -- (b#1-\i); \fi
    }}
  \baris{0}{fondasi ML dan DL}{ch:data,ch:stat,ch:teori,ch:optimasi,ch:dlbasic,ch:praktik,ch:lstm,ch:nlp}
  \baris{-1}{sains dan simulasi}{ch:optimasi,ch:fiskom,ch:fluida,ch:numerik,ch:pinn,ch:operator,ch:penemuan,ch:hibrida,ch:rcfd}
  \baris{-2}{ilmu hayati}{ch:data,ch:unsup,ch:gdl,ch:genomik,ch:hayati,ch:protein,ch:medis}
  \baris{-3}{sinyal dan robot}{ch:prob,ch:sinyal,ch:kendali,ch:pervasif,ch:robot,ch:rl}
  \baris{-4}{AI dan masyarakat}{ch:xai,ch:rekomendasi,ch:manusia,ch:keselamatan,ch:terbaru}
\end{tikzpicture}
\caption{Lima jalur baca. Angka adalah nomor bab. Jalur terbawah hampir tidak membutuhkan matematika.}
\label{fig:jalur}
\end{figure}

\section{Mengapa simulasi}

Simulasi fisika seperti \istilah{computational fluid dynamics} (CFD) menyelesaikan persamaan
diferensial parsial di jutaan sel. Hasilnya akurat, tetapi mahal: satu kasus bisa berjalan
berhari-hari di klaster komputer. Machine learning menawarkan jalan pintas. Ia belajar dari
hasil simulasi sebelumnya, lalu menebak hasil baru dalam hitungan detik. Masalahnya, model yang
hanya belajar dari data tidak tahu apa-apa tentang kekekalan massa atau energi. Ia bisa
menghasilkan aliran yang terlihat meyakinkan tetapi melanggar fisika.

Saya suka membayangkan perbedaan ini seperti dua cara memasak. Juru masak pertama mengikuti
resep secara persis: takaran, suhu, waktu. Hasilnya bisa diulang, tetapi lambat, karena setiap
langkah harus ditimbang. Juru masak kedua memasak dengan lidah. Ia mencicipi, mengingat rasa
masakan sebelumnya, lalu menebak. Cepat, tetapi kadang meleset jauh. Juru masak terbaik memakai
keduanya: resep menjaga agar masakan tidak gagal total, lidah membuatnya cepat dan luwes. Solver
CFD adalah resep, machine learning adalah lidah, dan \emph{physics-informed machine learning}
adalah usaha memakai keduanya sekaligus.

Di antara dua ujung itu terbentang spektrum yang lebar. \Cref{fig:spektrum} menempatkan
pendekatan-pendekatan yang akan kita temui di sepanjang buku.

\begin{figure}[ht]
\centering
\begin{tikzpicture}[x=1cm,y=1cm]
  \shade[left color=biru!25,right color=oranye!25] (0,0) rectangle (13,0.5);
  \node[label,anchor=north] at (0.2,-0.05) {\shortstack[l]{hanya fisika\\(solver klasik)}};
  \node[label,anchor=north] at (12.8,-0.05) {\shortstack[r]{hanya data\\(black box)}};
  \foreach \x/\t/\b in {
     1.2/{CFD, FV}/{Bab \ref{ch:numerik}},
     3.0/{rCFD}/{Bab \ref{ch:rcfd}},
     4.6/{model hibrida,\\ closure}/{Bab \ref{ch:hibrida}},
     6.3/{PINN}/{Bab \ref{ch:pinn}},
     7.9/{SINDy, ROM}/{Bab \ref{ch:penemuan}},
     9.6/{neural operator,\\ GNN simulator}/{Bab \ref{ch:operator}},
     11.6/{model generatif\\ murni}/{Bab \ref{ch:arsitektur}}}{
    \fill[abu] (\x,0.5) circle (0.06);
    \node[font=\sffamily\scriptsize,align=center,anchor=south] at (\x,0.62) {\t};
    \node[font=\sffamily\scriptsize,text=oranye!80!black,anchor=north] at (\x,-0.9) {\b};
  }
\end{tikzpicture}
\caption{Spektrum pendekatan simulasi, dari solver berbasis persamaan sampai model berbasis data.
Semakin ke kanan, semakin cepat, tetapi semakin sedikit jaminan fisikanya.}
\label{fig:spektrum}
\end{figure}

Literatur memakai beberapa nama untuk wilayah ini: \emph{scientific machine learning},
\emph{physics-informed machine learning} \citep{karniadakis2021}, atau \emph{physics-based deep
learning} \citep{thuerey2021}. Tinjauan yang khusus membahas mekanika fluida ditulis oleh
\citet{brunton2020}, \citet{vinuesa2022}, dan \citet{lino2023}. Di buku ini saya memakai istilah
yang lebih sederhana: AI untuk simulasi.

\section{Tiga pintu masuk fisika}

\citet{karniadakis2021} membedakan tiga cara memasukkan pengetahuan fisika ke dalam model
pembelajaran. Pembagian ini akan muncul berulang kali, jadi ada gunanya dikenal sejak awal.

Pintu pertama adalah data. Kalau sebuah model dilatih dengan cukup banyak hasil simulasi yang
benar, ia ikut meniru pola fisika yang terkandung di dalamnya. Inilah yang disebut
\istilah{observational bias}, dan contohnya adalah neural operator yang dilatih dari ribuan solusi
persamaan (\cref{ch:operator}). Pintu kedua adalah arsitektur. Model bisa dirancang agar
otomatis mematuhi sifat tertentu, misalnya tetap berlaku ketika sistemnya diputar atau digeser.
Konvolusi, yang dengan sendirinya menghormati pergeseran (\cref{ch:dlbasic}), dan jaringan yang
menghormati rotasi (\cref{ch:gdl}) termasuk di sini. Ini \istilah{inductive bias}. Pintu ketiga
adalah fungsi loss. Kita menambahkan suku yang menghukum setiap pelanggaran persamaan fisika,
sehingga model belajar mematuhinya selama pelatihan. Ini \istilah{learning bias}, dan contoh
paling terkenalnya adalah PINN (\cref{ch:pinn}).

Setiap kali bertemu model baru di bab-bab berikutnya, ada tiga pertanyaan yang layak diajukan.
Lewat pintu mana fisika masuk: data, arsitektur, atau loss? Apa yang terjadi kalau model
dipakai di luar jangkauan data latihnya, misalnya untuk waktu yang lebih panjang atau geometri
yang berbeda? Dan berapa biaya melatihnya dibanding biaya menjalankan solver klasik? Tiga
pertanyaan ini akan menemani kita sampai akhir.

\section{Notasi}

Buku ini memakai notasi yang umum di kuliah. Vektor dicetak tebal ($\vx$), matriks huruf
besar tebal ($\mW$), skalar huruf miring ($x$). Data latih ditulis $\{(\vx_n,y_n)\}_{n=1}^N$,
dan parameter model ditulis $\vw$ atau $\vtheta$. Ekspektasi ditulis $\E[\cdot]$, distribusi
normal $\N(\mu,\sigma^2)$, dan divergensi Kullback--Leibler $\KL(q\,\|\,p)$. Pada bab-bab
fluida, $\vu$ berarti kecepatan, $p$ tekanan, $\rho$ massa jenis, $\nu$ viskositas kinematik,
dan $T$ temperatur.
